\subsection{LOWER CENTRAL SERIES OF A FREE GROUP}

We shall prove the two following theorems:

THEOREM 2. Suppose that in the ring K the relation n.1 = 0 implies n = 0 for every integer n. Let r be a mapping of X into \(\hat{\mathrm{A}}(\mathrm{X})\) such that \(\omega(r(x)) \geqslant 2\) for \(x \in X\) and let g be the homomorphism of \(\mathrm{F}(\mathrm{X})\) into the Magnus group \(\Gamma(\mathrm{X})\) such that \(g(x) = 1 + x + r(x)\) for \(x \in X\) . For all \(n \geqslant 1\) , \(\mathrm{C}^n\mathrm{F}(\mathrm{X})\) is the inverse image under g of the subgroup \(1 + \hat{\mathrm{A}}_n(\mathrm{X})\) of \(\Gamma(\mathrm{X})\) .

THEOREM 3. For all \(x \in \mathbf{X}\) , let \(c(x)\) be the canonical image of \(x\) in \(\mathrm{F}(\mathbf{X}) / (\mathrm{F}(\mathbf{X}), \mathrm{F}(\mathbf{X}))\) . Let \(\mathfrak{g}\) be the graded Lie \(\mathbf{Z}\) -algebra associated with the filtration \((\mathrm{C}^{\mathrm{n}}\mathrm{F}(\mathbf{X}))_{n \geqslant 1}\) of \(\mathrm{F}(\mathbf{X})\) (§ 4, no. 6). The unique homomorphism of the free Lie \(\mathbf{Z}\) -algebra \(\mathrm{L}_{\mathbf{Z}}(\mathbf{X})\) into \(\mathfrak{g}\) which extends \(c\) is an isomorphism.

Loosely speaking, the graded Lie Z-algebra associated with the free group \(F(X)\) (with the lower central series) is the free Lie Z-algebra \(L_{Z}(X)\) .

We write \(\mathrm{F}(\mathbf{X}) = \mathrm{F}\) , \(\Gamma(\mathbf{X}) = \Gamma\) , \(\hat{\mathrm{A}}(\mathbf{X}) = \hat{\mathrm{A}}\) , \(\hat{\mathrm{A}}_{\mathbf{Z}}(\mathbf{X}) = \hat{\mathrm{A}}_{\mathbf{Z}}\) , \(\mathrm{C}^n\mathrm{F}(\mathbf{X}) = \mathrm{C}^n\) , \(\Gamma_n = 1 + \hat{\mathrm{A}}_n(\mathbf{X})\) and let \(\alpha: \mathrm{L}_{\mathbf{Z}}(\mathbf{X}) \to \mathfrak{g}\) be the homomorphism introduced in the statement of Theorem 3.

\textbf{(A) Preliminary reductions.}

Let \(\gamma\) denote the homomorphism of F into \(\Gamma\) defined by \(\gamma(x) = 1 + x\) for \(x \in \mathbf{X}\) . By Lemma 4 there exists an automorphism \(\sigma\) of the algebra \(\hat{\mathbf{A}}\) compatible with the filtration on \(\hat{\mathbf{A}}\) and such that \(\sigma(1 + x) = g(x)\) for all \(x \in \mathbf{X}\) ; then \(\sigma(\Gamma_n) = \Gamma_n\) for all \(n\) . As the homomorphisms \(g\) and \(\sigma \circ \gamma\) of F into \(\Gamma\) coincide on \(\mathbf{X}\) , \(g = \sigma \circ \gamma\) and hence \(g^{-1}(\Gamma_n) = \gamma^{-1}(\Gamma_n)\) . Under the hypotheses of Theorem 2, \(\mathbf{Z}\) can be identified with a subring of \(\mathbf{K}\) ; the Magnus algebra \(\hat{\mathbf{A}}_{\mathbf{Z}}\) is therefore identified with a subring of \(\hat{\mathbf{A}}\) and the filtration on \(\hat{\mathbf{A}}_{\mathbf{Z}}\) is induced by that on \(\hat{\mathbf{A}}\) . As \(\gamma\) maps F into \(\hat{\mathbf{A}}_{\mathbf{Z}}\) , we see that it suffices to prove Theorems 2 and 3 under the supplementary hypotheses \(\mathbf{K} = \mathbf{Z}\) , \(r = 0\) and hence \(g = \gamma\) , hypotheses which we shall henceforth make.

\textbf{(B) Surjectivity of \(\alpha\).}

As \(\mathbf{X}\) generates the group \(\mathbf{F} = \mathbf{C}^1\) , the set \(c(\mathbf{X})\) generates the \(\mathbf{Z}\) -module \(\mathfrak{g}^1 = \mathbf{C}^1 / \mathbf{C}^2\) . But \(\mathfrak{g}^1\) generates the Lie \(\mathbf{Z}\) -algebra \(\mathfrak{g}\) (§ 4, no. 6, Proposition 5) and hence \(c(\mathbf{X})\) generates \(\mathfrak{g}\) , which proves that \(\alpha\) is surjective.

\begin{enumerate}
\def\labelenumi{(\Alph{enumi})}
\setcounter{enumi}{2}
\tightlist
\item
  We identify the graded algebra \(\mathrm{gr}(\hat{\mathbf{A}})\) with A(X) under the canonical isomorphisms \(\mathrm{A}^n (\mathbf{X})\to \hat{\mathrm{A}}_n / \hat{\mathrm{A}}_{n + 1}\) . For every integer \(n\geqslant 1\) , we write \(\mathrm{F^n} = \gamma^{-1}(\Gamma_n)\) ; we know (§ 4, no. 5) that \((\mathbf{F}^n)_{n\geqslant 1}\) is an integral central filtration on F. Let \(\mathfrak{g}'\) denote the associated graded Lie Z-algebra (§ 4, no. 4). Let \(f\) be the Lie algebra homomorphism of \(\mathfrak{g}'\) into A(X) associated with \(\gamma\) (§ 4, no. 5, Proposition 3). Now \(\mathrm{C^n}\subset \mathrm{F^n}\) for every integer \(n\geqslant 1\) (§ 4, no. 6, Proposition 4) and hence there is a canonical homomorphism \(\varepsilon\) of \(\mathfrak{g} = \bigoplus_{n\geqslant 1}\mathrm{C}^n /\mathrm{C}^{n + 1}\) into \(\mathfrak{g}' = \bigoplus_{n\geqslant 1}\mathrm{F^n / F^{n + 1}}\)
\end{enumerate}

\[
\mathrm{L} _ {\mathbf {Z}} (\mathrm{X}) \stackrel {{\alpha}} {{\longrightarrow}} \mathfrak {g} \stackrel {{\varepsilon}} {{\longrightarrow}} \mathfrak {g} ^ {\prime} \stackrel {{f}} {{\longrightarrow}} \mathrm{A} (\mathrm{X}).
\]

We write \(\beta = f\circ \varepsilon\) ; we give \(\beta\) explicitly as follows: if \(u\) is the class modulo \(C^{n + 1}\) of an element \(w\) of \(C^n\) , then \(\gamma (w) - 1\) is of order \(\geqslant n\) in \(\hat{\mathbf{A}}\) and \(\beta (u)\) is the homogeneous component of \(\gamma (w) - 1\) of degree \(n\) . In particular,

\[
\beta (c (x)) = x \quad \text { for   all } x \in \mathbf {X}.\tag{3}
\]
% semantic-fragments: legacy-input-seam
\relax{}
\textbf{(D) Proof of Theorems 2 and 3.}

The Lie algebra homomorphism \(\beta \circ \alpha: \mathrm{L}_{\mathbf{Z}}(\mathrm{X}) \to \mathrm{A}(\mathrm{X})\) restricted to \(\mathbf{X}\) is the identity by (3) and hence is the canonical injection (§ 3, no. 1). Therefore \(\alpha\) is injective and hence bijective by (B); this proves Theorem 3. As \(\beta \circ \alpha = f \circ \varepsilon \circ \alpha\) is injective and \(\alpha\) is bijective, \(\varepsilon\) is injective. For all \(n \geqslant 1\) ,

\[
\varepsilon_ {n} \colon \mathrm{C} ^ {n} / \mathrm{C} ^ {n + 1} \rightarrow \mathrm{F} ^ {n} / \mathrm{F} ^ {n + 1}
\]

is injective and hence

\[
\mathrm{C} ^ {n} \cap \mathrm{F} ^ {n + 1} = \mathrm{C} ^ {n + 1}.
\]

\(C^{1} = F = F^{1}\) ; if \(C^{n} = F^{n}\) , then \(C^{n} \cap F^{n+1} = F^{n+1}\) whence \(C^{n+1} = F^{n+1}\) which proves Theorem 2 by induction on \(n \geqslant 1\) .

COROLLARY.

\[
\bigcap_ {n \geqslant 1} \mathrm{C} ^ {n} \mathrm{F} (\mathrm{X}) = \{e \}.
\]

Applying Theorem 2 with K = Z and r = 0,

\[
\bigcap_ {n \geqslant 1} \mathrm{C} ^ {n} \mathrm{F} (\mathrm{X}) = \bigcap_ {n \geqslant 1} \bar {g} (1 + \hat {\mathrm{A}} _ {n} (\mathrm{X})) = \bar {g} \left(\bigcap_ {n \geqslant 1} (1 + \hat {\mathrm{A}} _ {n} (\mathrm{X}))\right) = \bar {g} (1) = \{e \}.
\]

Remark. Let H be a Hall set relative to X (§2, no. 10). Let M be the magma defined by the law of composition \((x,y)\mapsto (x,y) = x^{-1}y^{-1}xy\) on F(X) and let \(\phi\) be the homomorphism of M(X) into M whose restriction to X is the identity. The elements of \(\phi (\mathrm{H})\) are called the basic commutators of F(X) associated with the Hall set H. For every integer \(n\geqslant 1\) , let \(\mathrm{H}_{\mathfrak{n}}\) be the subset of H consisting of the elements of length \(n\) ; we know (§2, no. 11, Theorem 1) that the canonical mapping of \(\mathrm{H}_{\mathfrak{n}}\) into \(\mathbf{L}_{\mathbf{Z}}(\mathbf{X})\) is a basis of the Abelian group \(\mathbf{L}_{\mathbf{Z}}^{n} (\mathbf{X})\) . Moreover, \(\phi (\mathrm{H_n})\subset \mathrm{C^n}\) ; for all \(m\in \mathrm{H}_{\mathfrak{n}}\) , let \(\phi_{\mathfrak{n}}(m)\) denote the class mod. \(\mathrm{C}^{n + 1}\) of \(\phi (m)\in \mathrm{C}^n\) . Theorem 3 then shows that \(\phi_{\mathfrak{n}}\) is a bijection of \(\mathrm{H}_{\mathfrak{n}}\) onto a basis of the Abelian group \(\mathrm{C}^n /\mathrm{C}^{n + 1}\) . It follows immediately that, for all \(w\in F(\mathbf{X})\) and all \(i\geqslant 1\) , there exists a unique element \(\alpha_{i}\) of \(\mathbf{Z}^{(\mathrm{H}_{i})}\) such that, for \(n\geqslant 1\) ,

\[
w = \prod_ {i = 1} ^ {n} \prod_ {m \in H _ {i}} \phi (m) ^ {\alpha_ {i} (m)} \mod C ^ {n + 1},\tag{4}
\]

where the product is calculated according to the total ordering given on H.

Example. Suppose that X is a set with two elements x, y and let \(H_{1} = \{x, y\}\) , \(H_{2} = \{xy\}\) . Every element w of F(X) can therefore be written

\[
w \equiv x ^ {a} y ^ {b} (x, y) ^ {c} \bmod . \mathrm{C} ^ {3} \quad \text { with } a, b, c \text { in } \mathbf {Z}.
\]

For \(w = (xy)^n\) , \(a = b = n\) and \(c = n(1 - n)/2\) (cf.~Exercise 9), whence

\[
(x y) ^ {n} \equiv x ^ {n} y ^ {n} (x, y) ^ {n (1 - n) / 2} \mod C ^ {3}.
\]

